\centerline{\bf \Huge Экзаменнационная работа по алгебре}
\centerline{\bf \Huge Вариант Q}

\begin{flushright}
        \scriptsize{Думали сессия закончилась?)}
\end{flushright}

\problem{1} 
Докажите, что произведение любых четырёх последовательных целых чисел при увеличении на 1 даёт точный квадрат.

\problem{2}
Назовём натуральное число \textit{почти квадратом}, если оно равно произведению двух последовательных натуральных чисел. Докажите, что каждый почти квадрат можно представить в виде частного двух почти квадратов.

\problem{3}
Известно, что $a+b=5, ab=2.$ Найдите:\\
а) $a^2 + b^2;$\\
б) $a^3 + b^3;$\\
в) $a^5 + b^5$

\problem{4}
Сто действительных чисел таковы, что каждое из них равно квадрату суммы всех остальных. Найдите сумму этих ста чисел.

\problem{5}
Вычислите значение выражения:

а) \((2 + 3)(2^2 + 3^2)(2^4 + 3^4) \ldots (2^{512} + 3^{512})\);

б) $\frac{2^3 - 1}{2^3 + 1} \cdot \frac{3^3 - 1}{3^3 + 1} \cdot \ldots \cdot \frac{2023^3 - 1}{2023^3 + 1} \cdot \frac{2023^3 + 1}{2023^3 + 1}$;

в) \(\frac{3^4 + 4}{3^4 - 4} + \frac{7^4 + 4}{5^4 - 4} + \frac{11^4 + 4}{9^4 - 4} + \ldots + \frac{2023^4 + 4}{2021^4 + 4}\).

\problem{6}
Пусть $a$ и $b$ $-$ целые числа. Докажите, что если $a^2+9ab+b^2$ делится на
$11$, то и $a^2-b^2$ делится на $11$.

\problem{7}
Найдите все такие натуральные $x, y$ и простые $p,$ что выполняется
равенство
\[x^3 + 3xy(x+y) + 2y^3 = p^2 \]

\problem{8}
При каких натуральных $n$ в выражении
\[\pm 1 \pm 2 \pm 3 \pm \ldots \pm n^2\]

можно так выбрать один из двух знаков $+$ или $ - $ , что в результате получится $0$?   

%https://problems.ru/view_problem_details_new.php?id=61454
